\documentclass[12pt]{article}
\usepackage[utf8]{inputenc}
\usepackage[T1]{fontenc}
\usepackage{hyperref}
\title{Plains Indian Sign Language Language Resources for K-12 Curriculum Development}
\author{Compiled by Claude Sonnet 5.5 (claude-sonnet-5-5) for Andr\'as Kornai}
\date{}
\begin{document}
\maketitle

\section{Introduction}
Plains Indian Sign Language (also Plains Sign Talk, Hand Talk, PISL/PSL; ISO 639-3 psd) is an endangered intertribal sign language of the Great Plains of the United States and Canada, now used as a first language only by elderly signers. Web searches found no ready-made school curriculum, primer or free online dictionary. What exists is mainly scholarly documentation and historical dictionaries in English.

\section{Primary Educational Resources}
No K-12 textbooks, primers or classroom materials in Plains Indian Sign Language were found in search results. Any curriculum would have to be built from the documentation below, in partnership with tribal communities.

\section{Linguistic and Reference Materials}
\begin{itemize}
\item RAPID: Plains Indian Sign Language Documentation, Linguistic Database, and Dictionary. An NSF-funded University of Arizona project (lead investigator Melanie Cody) associated with Kiowa, Northern Cheyenne, Peigan, Pawnee, Lakota and Ute communities. \url{https://www.nptao.arizona.edu/rapid-plains-indian-sign-language-documentation-linguistic-database-dictionary-0} and \url{https://naair.arizona.edu/rapid-plains-indian-sign-language-documentation-linguistic-database-dictionary-0}
\item Aboriginal Sign Languages of the Americas and Australia, Volumes I and II: a one-directional dictionary organised by English headword with written descriptions of signs (no audio, images or example phrases). Described by the UBC Relational Lexicography knowledgebase. \url{https://knowledgebase.arts.ubc.ca/aboriginal-sign-languages-of-the-americas-and-australia-volumes-i-and-ii}
\item UBC knowledgebase entry. \url{https://knowledgebase.arts.ubc.ca/?p=5773}
\item W.\ Tomkins, \emph{Universal Indian Sign Language of the Plains Indians of North America} (1927; Dover reissue 1969), over 525 signs; a 1927 copy is listed in a university repository. \url{https://repository.stou.ac.th/handle/6625047444/2525}
\end{itemize}

\section{Religious and Community Resources}
None found in search results.

\section{Digital and Audio Resources}
No free searchable online dictionary, video corpus or app was found in search results. The Encyclopedia.com entry gives a general overview. \url{https://www.encyclopedia.com/history/dictionaries-thesauruses-pictures-and-press-releases/indian-sign-language}

\section{Curriculum Development Context}
The language is endangered and held by elderly signers, so any curriculum for deaf and hearing learners depends on community-led documentation and permission. The RAPID project is the most relevant current contact point; historical dictionaries are written, English-based and lack video.
\end{document}